\section{Dataset Specifications}\label{app:datasets}

\looseness=-1 This appendix documents what we take from each upstream benchmark, how we split it, what the agent sees at solve time, and how grading runs.
Pinned upstream commits make the partition and grader reproducible from a clean checkout.

\subsection{SWE-Bench Pro}

SWE-Bench Pro \citep{deng2025swebenchpro} is the long-horizon software-engineering benchmark in our suite.
Tasks resolve only when a multi-file patch passes the upstream test set, and failure modes are therefore concrete and traceable to the harness's repository conventions and build commands.

\noindent\textbf{Source.}
We load the test split of \texttt{ScaleAI/SWE-bench\_Pro} from Hugging Face.
Evaluator scripts and per-instance Docker images come from the official \texttt{scaleapi/SWE-bench\_Pro-os} repository, pinned to a fixed commit.\footnote{Commit \texttt{0c64e26f00b9}.}

\noindent\textbf{Split.}
Rows are ordered by the SHA-256 hash of \texttt{(seed, instance\_id)} with the seed fixed.
The first 100 form the training pool and the next 100 the held-out test pool, and the remaining rows are unused.
We report on the held-out test pool.

\looseness=-1\noindent\textbf{Solve interface.}
Each task materializes a prompt file describing the issue plus a fresh clone of the upstream repository at its base commit, mounted in the workspace.
The agent edits repository files in place, and no extra tools are injected.

\noindent\textbf{Grading.}
We extract the agent's patch by re-applying its workspace edits to a fresh checkout and taking \texttt{git diff --binary}.
If the workspace path fails, we fall back to scanning the agent's final message for fenced diff blocks.
Binary hunks and auto-generated paths (dependency, cache, and build directories) are stripped before scoring.
The patch is then applied inside the official per-instance Docker image distributed by the SWE-Bench Pro authors, and the official run and parser scripts are invoked.
A task passes iff every \texttt{FAIL\_TO\_PASS} and \texttt{PASS\_TO\_PASS} test resolves correctly.
The Docker wall-clock budget is one hour per task.

\subsection{Terminal-Bench~2}

Terminal-Bench~2 \citep{terminalbench2} contains executable command-line tasks. Failure is a missed reward, not a missed convention, so the harness mostly buys consistency in how the agent inspects state and chains commands.

\noindent\textbf{Source.}
We use the upstream Terminal-Bench~2 repository at a fixed commit\footnote{Commit \texttt{53ff2b87d621}.}, yielding 89 tasks.

\noindent\textbf{Split.}
Tasks are hash-ordered with the same seeded scheme as SWE-Bench Pro.
The first 30 form the training pool, the remaining 59 the held-out pool.

\noindent\textbf{Solve interface.}
Each task spins up a fresh Docker container whose image is declared by the task.
Containers carry a cleanup label and per-task CPU and memory limits.
The agent's host workspace is bind-mounted inside the container.
The solve prompt instructs the agent to author shell scripts on the host and execute them inside the container.
A wall-clock watchdog terminates the container when the task's declared agent timeout elapses.

\looseness=-1\noindent\textbf{Grading.}
We run the task's upstream test suite inside the still-running container under its declared verifier timeout.
The verifier writes a reward of 0 or 1 to a known path.
A task passes iff the reward equals 1.
We apply no difficulty filter, so held-out numbers span the upstream difficulty mix.

\noindent\textbf{Data integrity.}
The solve prompt instructs the agent not to read the test directory, the upstream solution, or the verifier log.
This is a prompt-level convention, not a sandbox restriction.
We document it for transparency, since a sufficiently adversarial agent could read the verifier and game the reward.
We rely on the held-out pool and Table~\ref{tab:main}'s cross-benchmark consistency to detect such behavior, and we have not observed it.

\subsection{GAIA-2}

\looseness=-1 GAIA-2 \citep{froger2026gaia2} differs from the two coding benchmarks in that the environment evolves independently of the agent.
The harness layer therefore has to encode behavior under partial observation and asynchronous events, not just stable build conventions.

\noindent\textbf{Source.}
We load the validation split of the GAIA-2 release on Hugging Face under the \texttt{mini} configuration, yielding 200 scenarios.
Each scenario ships an asynchronous event stream and an upstream write-action verifier.

\noindent\textbf{Split.}
Scenarios are hash-ordered with the same seeded scheme.
The first 100 form the training pool, and the remainder, 100, form the held-out pool.
We report on the held-out slice.

\looseness=-1\noindent\textbf{Solve interface.}
Each scenario's workspace contains a task prompt, a tool dispatcher, and a tool catalog.
The agent invokes tools through a single shell command per call.
The dispatcher relays each call to a sidecar process running the upstream environment, which advances simulated time and replays scheduled events.
The scenario's ground-truth state is held entirely by the sidecar, and the agent never reads it.

\looseness=-1\noindent\textbf{Grading.}
When the agent finishes, the sidecar invokes the scenario's upstream verifier.
The judge LLM is the same Codex \texttt{gpt-5.5} backbone we use elsewhere, and routing defaults to the same Azure Foundry endpoint.
A scenario passes iff the upstream verifier reports success.

\noindent\textbf{Environment modifications.}
Two changes to the upstream environment affect reported numbers and must be disclosed.
First, the upstream cap of one \texttt{send\_message\_to\_user} call per turn is raised to four, since the original cap penalizes verbose agents that would otherwise complete the task correctly.
Second, three optional judge-relaxation switches (event filtering, relaxed per-app UI judging, and trivial filesystem-read filtering) are available as ablation axes but are all disabled in the experiments reported in Table~\ref{tab:main}.
